% This file is part of GUFI, which is part of MarFS, which is released
% under the BSD license.
%
%
% Copyright (c) 2017, Los Alamos National Security (LANS), LLC
% All rights reserved.
%
% Redistribution and use in source and binary forms, with or without modification,
% are permitted provided that the following conditions are met:
%
% 1. Redistributions of source code must retain the above copyright notice, this
% list of conditions and the following disclaimer.
%
% 2. Redistributions in binary form must reproduce the above copyright notice,
% this list of conditions and the following disclaimer in the documentation and/or
% other materials provided with the distribution.
%
% 3. Neither the name of the copyright holder nor the names of its contributors
% may be used to endorse or promote products derived from this software without
% specific prior written permission.
%
% THIS SOFTWARE IS PROVIDED BY THE COPYRIGHT HOLDERS AND CONTRIBUTORS "AS IS" AND
% ANY EXPRESS OR IMPLIED WARRANTIES, INCLUDING, BUT NOT LIMITED TO, THE IMPLIED
% WARRANTIES OF MERCHANTABILITY AND FITNESS FOR A PARTICULAR PURPOSE ARE DISCLAIMED.
% IN NO EVENT SHALL THE COPYRIGHT HOLDER OR CONTRIBUTORS BE LIABLE FOR ANY DIRECT,
% INDIRECT, INCIDENTAL, SPECIAL, EXEMPLARY, OR CONSEQUENTIAL DAMAGES (INCLUDING,
% BUT NOT LIMITED TO, PROCUREMENT OF SUBSTITUTE GOODS OR SERVICES; LOSS OF USE,
% DATA, OR PROFITS; OR BUSINESS INTERRUPTION) HOWEVER CAUSED AND ON ANY THEORY OF
% LIABILITY, WHETHER IN CONTRACT, STRICT LIABILITY, OR TORT (INCLUDING NEGLIGENCE
% OR OTHERWISE) ARISING IN ANY WAY OUT OF THE USE OF THIS SOFTWARE, EVEN IF
% ADVISED OF THE POSSIBILITY OF SUCH DAMAGE.
%
%
% From Los Alamos National Security, LLC:
% LA-CC-15-039
%
% Copyright (c) 2017, Los Alamos National Security, LLC All rights reserved.
% Copyright 2017. Los Alamos National Security, LLC. This software was produced
% under U.S. Government contract DE-AC52-06NA25396 for Los Alamos National
% Laboratory (LANL), which is operated by Los Alamos National Security, LLC for
% the U.S. Department of Energy. The U.S. Government has rights to use,
% reproduce, and distribute this software.  NEITHER THE GOVERNMENT NOR LOS
% ALAMOS NATIONAL SECURITY, LLC MAKES ANY WARRANTY, EXPRESS OR IMPLIED, OR
% ASSUMES ANY LIABILITY FOR THE USE OF THIS SOFTWARE.  If software is
% modified to produce derivative works, such modified software should be
% clearly marked, so as not to confuse it with the version available from
% LANL.
%
% THIS SOFTWARE IS PROVIDED BY LOS ALAMOS NATIONAL SECURITY, LLC AND CONTRIBUTORS
% "AS IS" AND ANY EXPRESS OR IMPLIED WARRANTIES, INCLUDING, BUT NOT LIMITED TO,
% THE IMPLIED WARRANTIES OF MERCHANTABILITY AND FITNESS FOR A PARTICULAR PURPOSE
% ARE DISCLAIMED. IN NO EVENT SHALL LOS ALAMOS NATIONAL SECURITY, LLC OR
% CONTRIBUTORS BE LIABLE FOR ANY DIRECT, INDIRECT, INCIDENTAL, SPECIAL,
% EXEMPLARY, OR CONSEQUENTIAL DAMAGES (INCLUDING, BUT NOT LIMITED TO, PROCUREMENT
% OF SUBSTITUTE GOODS OR SERVICES; LOSS OF USE, DATA, OR PROFITS; OR BUSINESS
% INTERRUPTION) HOWEVER CAUSED AND ON ANY THEORY OF LIABILITY, WHETHER IN
% CONTRACT, STRICT LIABILITY, OR TORT (INCLUDING NEGLIGENCE OR OTHERWISE) ARISING
% IN ANY WAY OUT OF THE USE OF THIS SOFTWARE, EVEN IF ADVISED OF THE POSSIBILITY
% OF SUCH DAMAGE.



\subsection{Virtual Tables}
A series of
\href{https://www.sqlite.org/vtab.html}{\sqlite~Virtual~Tables} have
been developed for GUFI: \gufivtstar, \gufivt, and \runvt. These
virtual tables are available in \gufisqlite, and as
\href{https://sqlite.org/loadext.html}{runtime loadable
extensions}\footnote{\gufivtstar and \gufivt are in the same
binary.}\textsuperscript{,}\footnote{Remember to make \gufiquery
available via \texttt{PATH}.}. This allows for indexes to be used by
\sqlite libraries in any language that can load \sqlite extensions,
such as the C, Python, R, or Ruby libraries. The \sqlite shell can
also be used to load these extensions.

\subsubsection{\gufivtstar}
\gufivtstar refers to the following collection of virtual tables:
\begin{itemize}
\item \gufivttreesummary
\item \gufivtsummary
\item \gufivtentries
\item \gufivtpentries
\item \gufivtvrsummary
\item \gufivtvrpentries
\end{itemize}

They each wrap \gufiquery with fixed SQL queries. Tree traversal
behavior maybe customized via positional arguments. Unused trailing
arguments may be left out. \texttt{NULL} may be used to provide values
to unused arguments that come before used arguments.

\begin{table}[H]
  \centering
  \caption{\gufivtstar argument order}
  \begin{tabularx}{\textwidth}{| l | c | X |}
    \hline
    Argument & Type & Notes \\
    \hline
    index path & String & \\
    \hline
    \# of threads & Positive integer & \\
    \hline
    min\_level & Non-negative integer & \\
    \hline
    path\_list & String & File path \\
    \hline
    verbose & 0 or 1 & Print command before running it \\
    \hline
    remote\_cmd & String & Command to forward this run of \gufiquery to a remote (e.g. \texttt{ssh}) \\
    \hline
    remote\_args & String & Space separated arguments to pass to
    \texttt{remote\_cmd} \\
    \hline
  \end{tabularx}
\end{table}

\gufivtstar virtual table names are directly usable as
\href{https://www.sqlite.org/vtab.html#tabfunc2}{table-valued~functions}. The
result of each \gufivtstar is presented as a single \sqlite virtual table with a
fixed schema. \\

Example Usage:
\begin{verbatim}
(
    echo ".load /path/to/gufi_vt" # do not provide extension
    echo "SELECT * FROM gufi_vt_pentries(`/path/to/index', 2);"
) | sqlite3
\end{verbatim}

\subsubsection{\gufivt}
\indent \gufivt is a more generic version of \gufivtstar that allows
for arbitrary queries to be run on an index. However, because the
resultant table schema is not known beforehand, a virtual table must
be set up before being used.

The \gufivt arguments are mostly the same as those of \gufivtstar. The
difference are:

\begin{itemize}
\item Arguments are \texttt{key=value} pairs
\item Order of arguments with different keys does not matter
\item Unless otherwise noted, arguments with duplicate keys results in
  the rightmost argument being used
\item There are a few arguments that \gufivtstar does not have
\end{itemize}

The ordering of the following table is different from the \gufivtstar
argument table in order to better reflect the command that is run.

\begin{longtable}{| l | c | p{0.35\linewidth} |}
  \hline
  Key & Value Type & Notes \\
  \hline
  remote\_cmd & String & Command to forward this run of \gufiquery
  to a remote (e.g. \texttt{ssh}) \\
  \hline
  remote\_arg & String & Single \texttt{argv[i]} passed to
  \texttt{remote\_cmd} before \gufiquery (e.g. \texttt{user@remote})
  \\
  & & \\
  & & Multiple are allowed \\
  \hline
  \# of threads & Positive Integer & \\
  \hline
  min\_level & Non-negative Integer & \\
  \hline
  dir\_match\_uid & UID & Not available in \gufivtstar \\
  \hline
  dir\_match\_gid & GID & Not available in \gufivtstar \\
  \hline
  setup\_res\_col\_type & SQL & Set up temporary environment to make
  namespaces/tables/columns available for getting result column
  types \\
  \hline
  I, T, S, E, K, J, G, F & SQL & Forwarded to \gufiquery \\
  \hline
  path\_list & File Path & \\
  \hline
  p & Directory Path & Source path for use with \texttt{spath()} \\
  \hline
  plugin & Formatted String & \textless
  entrypoint\textgreater:\textless gufi\_query plugin library
  path\textgreater \\
  & & \\
  & & Multiple are allowed \\
  & & \\
  & & Not available in \gufivtstar \\
  \hline
  index & Directory Path & Arguments without keys are also treated
  as the index path \\
  & & \\
  & & Multiple are allowed \\
  \hline
  verbose or VERBOSE & 0 or 1 & Print command before running it \\
  \hline
\end{longtable}

Example Usage:
\begin{verbatim}
(
    echo ".load /path/to/gufi_vt" # do not provide extension
    echo "CREATE VIRTUAL TABLE temp.output" \
         "USING gufi_vt(index=/path/to/index, threads=2," \
         "E=`SELECT name, size FROM pentries');"
    echo "SELECT * FROM output;"
) | sqlite3
\end{verbatim}

\subsubsection{\runvt}
\runvt is the generic virtual table that runs input commands with
\texttt{popen(3)} and converts each line of the output into a row of a
\sqlite table.

\runvt is available in \gufiquery, \gufisqlite, and as a runtime
loadable extension.
\\\\
Example Usage:
\begin{verbatim}
(
    echo ".load /path/to/run_vt" # do not provide extension
    echo "CREATE VIRTUAL TABLE temp.output " \
         "USING run_vt(cmd=`cat file.csv', " \
         "d=`,', cols=`name, size');"
    echo "SELECT * FROM output;"
) | sqlite3
\end{verbatim} 
